\chapter{Nhận diện món ăn Việt Nam qua ảnh}

\section{Bài toán và hướng tiếp cận}

Một ứng dụng ẩm thực hoặc du lịch cho phép người dùng chụp ảnh món ăn và tự động cho biết tên món (phở, bún bò Huế,
bánh xèo, cơm tấm, ...). Bài tập này xây dựng mô hình nhận một ảnh và dự đoán món ăn trong ảnh, trong 30 món ăn Việt
Nam phổ biến. Vì ứng dụng có thể hiện vài gợi ý cùng lúc, ngoài accuracy và macro-F1, nhóm còn báo cáo top-3 và top-5
accuracy: tỉ lệ ảnh có món đúng nằm trong 3 hoặc 5 dự đoán đầu tiên.

Đây là bài toán phân loại ảnh có giám sát với 30 lớp. Hướng chính là học sâu với một mạng đã huấn luyện sẵn: ResNet50
\cite{img-resnet} được dùng theo hai cách mà đề yêu cầu, trích đặc trưng (giữ nguyên mạng) và fine-tune nhẹ (học lại
khối cuối), và được so sánh với một baseline dùng đặc trưng thủ công (histogram màu và SIFT). Cả ba mô hình có cùng một
kiểu đầu phân loại, nên khác biệt giữa chúng chỉ đến từ đặc trưng. Quy trình gồm: đưa dữ liệu về cấu trúc thư mục chuẩn
cho phân loại ảnh; kiểm tra chất lượng và phân tích khám phá; tiền xử lý và tăng cường dữ liệu; xây dựng ba mô hình;
đánh giá một lần trên tập test và phân tích lỗi. Notebook chạy được trên Google Colab với \textit{Runtime > Run all}
(cần runtime GPU).

\section{Bộ dữ liệu}

\begin{itemize}
  \item \textbf{Nguồn:} bộ dữ liệu 30VNFoods \cite{img-30vnfoods}, bản gốc được đăng trên Kaggle.
        Notebook tải bản sao công khai trên Hugging Face (\url{vuongnhathien/30VNFoods}), không cần đăng nhập.
  \item \textbf{Giấy phép:} dữ liệu được phát hành theo giấy phép CC BY-NC-SA 4.0, chỉ dùng cho mục đích phi thương
        mại và phải ghi nguồn. Nhóm không đưa ảnh lên repo.
  \item \textbf{Nội dung:} ảnh chụp 30 món ăn Việt Nam thu thập từ Internet. Mỗi ảnh có một nhãn là tên món, lưu bằng tên
        thư mục chứa ảnh (tên món viết không dấu).
  \item \textbf{Kết quả gốc:} mô hình tốt nhất của bài báo đạt accuracy 77.54\% và top-5 accuracy 96.07\% trên tập test.
\end{itemize}

Bộ dữ liệu có 25 136 ảnh của 30 món, nhiều hơn mức tối thiểu 5 000 ảnh và 3 lớp mà đề yêu cầu cho dữ liệu ảnh, và là
ảnh chụp thực tế chứ không phải bộ dữ liệu mẫu như MNIST hay CIFAR-10. Tác giả đã chia sẵn thành train (69.9\%), val
(10.0\%) và test (20.1\%); nhóm giữ nguyên cách chia này để kết quả so sánh được với bài báo gốc
(Bảng~\ref{tab:img-split}). Kích thước ảnh rất khác nhau: cạnh ngắn từ 83 đến 4 912 px, trung vị 700 px, và chỉ 9 ảnh
có cạnh ngắn dưới mức 128 px mà đề khuyến nghị.

\paragraph{Cấu trúc thư mục.} Mỗi ảnh chỉ có một nhãn, nên nhãn được lưu bằng thư mục chứa ảnh. Notebook đưa dữ liệu về
đúng cấu trúc mà Ultralytics YOLO dùng cho bài toán phân loại: một thư mục cho mỗi tập, bên trong là một thư mục con cho
mỗi món. File \texttt{labels.csv} liệt kê mọi ảnh với tập và nhãn của nó.

\begin{verbatim}
data/
  train/
    Banh beo/   1.jpg  10.jpg  ...
    Banh bot loc/
    ...         (30 thu muc, moi mon mot thu muc)
  val/          (cung cau truc)
  test/         (cung cau truc)
  labels.csv    (path, split, label)
\end{verbatim}

\begin{table}[H]\centering
\caption{Kích thước ba tập trước và sau khi làm sạch.}\label{tab:img-split}
\begin{tabular}{lrrrrr}
\toprule
Tập & Ban đầu & Ảnh nhỏ & Ảnh trùng & Còn lại & Số ảnh mỗi món \\
\midrule
train & 17 581 & 5 & 361 & 17 215 & 306 - 1 047 \\
val & 2 515 & 0 & 0 & 2 515 & 45 - 153 \\
test & 5 040 & 0 & 0 & 5 040 & 89 - 306 \\
\bottomrule
\end{tabular}
\end{table}

\section{Kiểm tra chất lượng dữ liệu}

Mỗi ảnh được mở một lần để ghi lại: file có đọc được không, định dạng thật theo nội dung file, chế độ màu, kích thước,
mã băm cảm nhận (pHash) và độ bão hòa, độ sáng trung bình. Bảng~\ref{tab:img-quality} tổng hợp kết quả.

\begin{table}[H]\centering
\caption{Các vấn đề chất lượng dữ liệu và cách xử lý.}\label{tab:img-quality}
\begin{tabular}{>{\raggedright\arraybackslash}p{4.4cm}>{\raggedright\arraybackslash}p{4.6cm}>{\raggedright\arraybackslash}p{5.2cm}}
\toprule
Vấn đề & Kết quả kiểm tra & Cách xử lý \\
\midrule
File lỗi & 0 file & nếu có thì bỏ khỏi train/val/test \\
Định dạng khác JPEG & 129 file (PNG, MPO, BMP) có đuôi \texttt{.jpg} & giữ lại, thư viện đọc ảnh theo nội dung file \\
Chế độ màu khác RGB & 0 ảnh & vẫn đổi sang RGB, cần cho ảnh mới \\
Cạnh ngắn dưới 128 px & 9 ảnh (5 ở train) & bỏ khỏi train; giữ ở val/test và phóng to \\
Kích thước khác nhau & cạnh ngắn từ 83 đến 4 912 px, trung vị 700 px & đưa về 224 x 224 \\
Ảnh trùng giữa train và val/test & 186 cặp & bỏ khỏi train \\
Ảnh trùng mang nhãn khác & 65 cặp & bỏ khỏi train; đánh dấu ở val/test \\
Ảnh trùng cùng nhãn trong train & 153 cặp & giữ một bản cho mỗi nhóm, bỏ các bản còn lại khỏi train \\
Lớp mất cân bằng & món nhiều ảnh nhất gấp 3.5 lần món ít nhất & chọn mô hình theo macro-F1 \\
\bottomrule
\end{tabular}
\end{table}

\reportfig{fig01_image_sizes.png}{Kích thước ảnh rất đa dạng nhưng phần lớn đủ lớn: cạnh ngắn trung vị 700 px, chỉ 9 ảnh
(0.04\%) dưới 128 px và 347 ảnh (1.4\%) dưới 224 px.}{fig:img-sizes}

Để tìm ảnh trùng, mỗi ảnh được tóm tắt bằng pHash, một mã 64 bit mô tả cấu trúc tổng thể của ảnh: hai bản của cùng một
ảnh (dù đã bị nén lại, thu nhỏ hay cắt nhẹ) có mã giống hoặc gần giống nhau. Nhóm so sánh toàn bộ 315 896 680 cặp ảnh
bằng phép XOR và đếm bit của numpy, và coi hai ảnh là trùng nếu mã của chúng khác nhau không quá 6 bit. Kết quả có 403
cặp trùng (239 cặp trùng hoàn toàn, 164 cặp gần trùng), liên quan đến 757 ảnh (3.0\% bộ dữ liệu).

\reportfig{fig02_duplicate_pairs.jpg}{Ví dụ các cặp ảnh trùng. Bốn cặp đầu (tên màu đỏ) là cùng một ảnh nhưng được xếp
vào hai món khác nhau, như bánh bèo và bánh bột lọc; hai cặp tiếp theo có bản sao ở cả train lẫn validation hoặc test;
hai cặp cuối là ảnh gần trùng trong train.}{fig:img-dups}

\section{Phân tích khám phá dữ liệu}

Phân tích khám phá chỉ dùng tập train, trừ Hình~\ref{fig:img-classes} so sánh ba tập.

\reportfig{fig03_class_counts.png}{Số ảnh của mỗi món theo tập. Bún bò Huế có nhiều ảnh train nhất (1 071 ảnh), bánh pía
ít nhất (310 ảnh), chênh 3.5 lần; ba tập giữ cùng tỉ lệ giữa các món.}{fig:img-classes}

Mất cân bằng chỉ ở mức vừa phải và món ít ảnh nhất vẫn có hơn 300 ảnh train, nên nhóm giữ mọi món, không lấy mẫu lại và
không dùng trọng số lớp, nhưng vẫn chọn mô hình theo macro-F1 để mọi món có trọng số như nhau.

\reportfig{fig04_sample_grid.jpg}{Bốn ảnh train ngẫu nhiên của mỗi món. Ảnh được chụp từ nhiều góc, nhiều kiểu bày đĩa và
điều kiện ánh sáng khác nhau. Nhiều món rất giống nhau, như các món nước (phở, bún bò Huế, hủ tiếu, bún riêu, bánh canh)
hay các loại bánh làm từ bột gạo (bánh bèo, bánh đúc, bánh bột lọc).}{fig:img-samples}

\reportfig{fig05_colour_stats.png}{Độ bão hòa và độ sáng trung bình của ảnh theo món. Trung vị độ bão hòa của các món chỉ
trải từ 0.33 (cháo lòng) đến 0.52 (bánh xèo), trong khi khoảng tứ phân vị bên trong một món trung bình đã rộng 0.16, nên
màu sắc tổng thể khó đủ để phân biệt 30 món.}{fig:img-colour}

\section{Tiền xử lý và tăng cường dữ liệu}

Validation và test không bị bỏ ảnh nào, để giữ nguyên cách chia chính thức; mọi thống kê học từ dữ liệu chỉ được fit
trên train.

\begin{enumerate}
  \item Làm sạch train: bỏ 5 ảnh có cạnh ngắn dưới 128 px, 182 ảnh trùng với validation hoặc test, 35 ảnh trùng mang nhãn
        khác và 144 ảnh lặp lại cùng nhãn; train còn 17 215 ảnh. Ảnh validation và test được đánh dấu nếu có bản gần trùng
        trong train trước khi làm sạch (123 ảnh test) hoặc có bản sao mang nhãn khác (37 ảnh test).
  \item Mã hóa nhãn bằng \texttt{LabelEncoder}: 30 tên món thành số nguyên 0-29.
  \item Thu nhỏ một lần và lưu lại: ResNet50 nhận ảnh 224 x 224 px, trong khi ảnh gốc có cạnh ngắn trung vị 700 px, và
        giải mã ảnh JPEG lớn ở mọi epoch rất chậm. Mọi ảnh được thu nhỏ để cạnh ngắn bằng 256 px (giữ tỉ lệ) và lưu vào
        \texttt{data/processed/} cùng cấu trúc thư mục; mọi mô hình đọc ảnh từ đây, ảnh gốc không bị thay đổi.
  \item Đưa ảnh vào mạng: đổi kích thước để cạnh ngắn bằng 256 px (ảnh nhỏ hơn được phóng to), cắt phần giữa
        224 x 224 px, chuẩn hóa từng kênh màu theo trung bình và độ lệch chuẩn của ImageNet. Ảnh luôn được cắt thành hình
        vuông nên không cần thêm viền.
  \item Chuẩn hóa vector đặc trưng trước Logistic Regression: chuẩn hóa L2 với đặc trưng ResNet50 (đưa mỗi vector về độ
        dài 1), \texttt{StandardScaler} với đặc trưng thủ công (vì hai phần của vector có thang đo khác nhau).
\end{enumerate}

\paragraph{Tăng cường dữ liệu.} Mỗi ảnh train được biến đổi ngẫu nhiên ở mỗi epoch, để mạng thấy nhiều phiên bản khác
nhau của cùng một ảnh, học những đặc điểm không đổi của món ăn thay vì học thuộc từng ảnh
(Bảng~\ref{tab:img-aug}). Tăng cường chỉ áp dụng cho train và chỉ khi trọng số của mạng được huấn luyện (fine-tune);
validation và test luôn dùng phép biến đổi cố định để điểm số không phụ thuộc vào yếu tố ngẫu nhiên. Ở cách trích đặc
trưng, mạng không được huấn luyện nên mỗi ảnh chỉ đi qua mạng một lần.

\begin{table}[H]\centering
\caption{Các phép tăng cường dữ liệu và mục tiêu.}\label{tab:img-aug}
\begin{tabular}{>{\raggedright\arraybackslash}p{3.4cm}>{\raggedright\arraybackslash}p{4.4cm}>{\raggedright\arraybackslash}p{6.4cm}}
\toprule
Phép biến đổi & Tham số & Mục tiêu \\
\midrule
Cắt khung ngẫu nhiên & 224 x 224, giữ 70-100\% diện tích & mô phỏng khung hình và khoảng cách chụp khác nhau; món ăn vẫn nằm trong khung \\
Lật ngang & xác suất 0.5 & ảnh lật ngang vẫn là cùng một món \\
Đổi màu nhẹ & độ sáng, tương phản $\pm$20\%; độ bão hòa $\pm$10\%; sắc màu $\pm$0.02 & mô phỏng ánh sáng khác nhau; sắc màu chỉ đổi rất ít vì màu là đặc điểm thật của món \\
\bottomrule
\end{tabular}
\end{table}

Không lật dọc hay xoay mạnh, vì ảnh món ăn thật hiếm khi bị như vậy.

\reportfig{fig06_augmentation.jpg}{Một ảnh train của món bánh xèo và sáu phiên bản tăng cường ngẫu nhiên: món ăn và màu
sắc của nó vẫn nhận ra được.}{fig:img-aug}

\section{Xây dựng mô hình}

Nhóm dùng ResNet50 với trọng số ImageNet của torchvision \cite{img-resnet}, một trong các mạng đề gợi ý. Mạng được coi
như một hàm có sẵn: nhận ảnh 224 x 224 và trả về một vector 2 048 chiều mô tả nội dung ảnh; lớp cuối biến vector này
thành xác suất của từng lớp. Ba mô hình được xây dựng (Bảng~\ref{tab:img-models}). Logistic Regression nhiều lớp là một
lớp tuyến tính cộng softmax, đúng bằng lớp cuối của ResNet50, nên cả ba mô hình có cùng một kiểu đầu phân loại và chỉ
khác nhau ở đặc trưng. Mô hình được huấn luyện trên train đã làm sạch và chọn theo macro-F1 trên validation; tập test
chỉ dùng một lần. Logistic Regression dùng tham số mặc định của scikit-learn.

\begin{table}[H]\centering
\caption{Ba mô hình được so sánh.}\label{tab:img-models}
\begin{tabular}{>{\raggedright\arraybackslash}p{3.6cm}>{\raggedright\arraybackslash}p{4.2cm}>{\raggedright\arraybackslash}p{3.2cm}>{\raggedright\arraybackslash}p{3.2cm}}
\toprule
Mô hình & Đặc trưng & Phần được huấn luyện & Đầu phân loại \\
\midrule
ResNet50 trích đặc trưng & vector 2 048 chiều của ResNet50 giữ nguyên & đầu phân loại & Logistic Regression \\
ResNet50 fine-tune nhẹ & ResNet50, khối cuối \texttt{layer4} học lại & \texttt{layer4} và đầu phân loại & tuyến tính 2048 $\rightarrow$ 30 + softmax \\
Baseline thủ công & histogram màu + SIFT (OpenCV) & đầu phân loại & Logistic Regression \\
\bottomrule
\end{tabular}
\end{table}

\subsection{ResNet50 làm bộ trích đặc trưng}

ResNet50 bỏ lớp phân loại cuối và giữ nguyên toàn bộ trọng số; mỗi ảnh đi qua mạng một lần và trở thành một vector
2 048 chiều, rồi Logistic Regression được huấn luyện trên vector của tập train. Chỉ huấn luyện đầu phân loại đã đạt
macro-F1 validation 0.694 và top-3 0.873, trong khi đoán ngẫu nhiên chỉ đạt khoảng 0.033.

\subsection{Fine-tune nhẹ ResNet50}

Bắt đầu từ cùng ResNet50, lớp cuối được thay bằng một lớp tuyến tính 30 lớp mới; mọi lớp trừ khối cuối \texttt{layer4}
và lớp mới bị đóng băng (15.0 triệu trên 23.6 triệu tham số được huấn luyện). Mạng được huấn luyện 3 epoch bằng Adam với
learning rate nhỏ ($10^{-4}$) để không phá vỡ trọng số đã học, batch 64, có tăng cường dữ liệu, và epoch có macro-F1
validation cao nhất được giữ lại.

\reportfig{fig07_finetune_curves.png}{Hàm mất mát train và kết quả validation qua các epoch. Macro-F1 validation tăng từ
0.767 lên 0.817 rồi 0.840; đường chấm là macro-F1 của cách trích đặc trưng (0.694).}{fig:img-ft}

Sau fine-tune, macro-F1 validation là 0.840, cao hơn cách trích đặc trưng 14.6 điểm; ngay sau epoch đầu tiên mô hình đã
vượt cách trích đặc trưng. Validation vẫn đang tăng ở epoch 3, nên huấn luyện thêm có thể còn cải thiện.

\subsection{Baseline: đặc trưng thủ công}

Để biết mạng học sẵn đóng góp bao nhiêu, cùng một Logistic Regression được huấn luyện trên đặc trưng thiết kế bằng tay,
tính trên ảnh 128 x 128 px:

\begin{itemize}
  \item \textbf{Histogram màu} trong không gian HSV (OpenCV, 8 x 8 x 8 = 512 chiều): trên đĩa có những màu nào.
  \item \textbf{SIFT} (OpenCV): tìm các điểm nổi bật và mô tả vùng quanh mỗi điểm bằng một vector 128 chiều về hướng của
        các cạnh. Vì mỗi ảnh có số điểm khác nhau, nhóm dùng \textit{túi từ thị giác}: \texttt{MiniBatchKMeans} của
        scikit-learn gom các vector SIFT của 2 000 ảnh train thành 256 cụm, và mỗi ảnh được mô tả bằng tỉ lệ điểm rơi vào
        từng cụm (256 chiều).
\end{itemize}

Baseline đạt macro-F1 validation 0.241, cao hơn đoán ngẫu nhiên nhưng thấp hơn ResNet50 trích đặc trưng gần 3 lần: màu
sắc và hình dạng cục bộ không đủ để phân biệt các món có màu và cách bày giống nhau. Theo cách chọn mô hình,
\textbf{ResNet50 fine-tune nhẹ} (macro-F1 validation 0.840) là mô hình cuối cùng.

\section{Đánh giá trên tập test}

\begin{table}[H]\centering
\caption{Kết quả trên tập test (5 040 ảnh). Hàng in đậm là mô hình cuối cùng, chọn theo macro-F1 validation.}\label{tab:img-test}
\small
\begin{tabular}{llrrrr}
\toprule
Mô hình & Đặc trưng & Accuracy & Macro-F1 & Top-3 & Top-5 \\
\midrule
Baseline thủ công & histogram màu + SIFT & 0.2522 & 0.2400 & 0.4544 & 0.5683 \\
ResNet50 trích đặc trưng & ResNet50 giữ nguyên & 0.6952 & 0.6945 & 0.8762 & 0.9306 \\
\textbf{ResNet50 fine-tune nhẹ} & \textbf{layer4 học lại} & \textbf{0.8363} & \textbf{0.8397} & \textbf{0.9470} & \textbf{0.9738} \\
\midrule
Bài báo gốc & & 0.7754 & & & 0.9607 \\
\bottomrule
\end{tabular}
\end{table}

Mô hình cuối cùng đạt accuracy 0.8363 và top-5 0.9738, cao hơn bài báo gốc (0.7754 và 0.9607). Thứ tự ba mô hình trên
test giống trên validation, và macro-F1 trên test gần như bằng trên validation (0.8397 so với 0.8403).

\reportfig{fig08_model_comparison.png}{Ba mô hình trên cùng tập test, cùng một kiểu đầu phân loại. Macro-F1 tăng từ 0.240
(đặc trưng thủ công) lên 0.695 (ResNet50 trích đặc trưng, +45.5 điểm) và 0.840 (fine-tune nhẹ, thêm 14.5
điểm).}{fig:img-features}

\reportfig{fig09_per_class_f1.png}{F1 của từng món trên test. F1 trải từ 0.63 (bánh canh) đến 0.97 (bánh mì); ba món khó
nhất (bánh canh, bánh đúc, hủ tiếu) đều thuộc nhóm món nước hoặc bánh làm từ bột gạo.}{fig:img-perclass}

\reportfig{fig10_confusion_matrix.png}{Ma trận nhầm lẫn chuẩn hóa theo hàng. Số màu đỏ là 15 ô ngoài đường chéo từ 5\%
trở lên, trên tổng số 870 ô: lỗi tập trung vào một số cặp món.}{fig:img-cm}

\subsection{Những món nào hay bị nhầm với nhau?}

\begin{table}[H]\centering
\caption{Năm cặp món bị nhầm nhiều nhất (tỉ lệ trên số ảnh test của món thực tế).}\label{tab:img-pairs}
\begin{tabular}{llrrr}
\toprule
Món thực tế & Bị đoán thành & Tỉ lệ & Số ảnh & Chiều ngược lại \\
\midrule
Bánh bèo & Bánh đúc & 12.4\% & 16 & 3.0\% \\
Bún riêu & Bún bò Huế & 10.4\% & 24 & 4.6\% \\
Bánh canh & Hủ tiếu & 10.4\% & 20 & 10.2\% \\
Hủ tiếu & Bánh canh & 10.2\% & 20 & 10.4\% \\
Bún mắm & Canh chua & 9.0\% & 14 & 2.4\% \\
\bottomrule
\end{tabular}
\end{table}

\reportfig{fig11_confused_pairs.jpg}{Ba lỗi của mỗi cặp nhầm nhiều nhất (đỏ) cạnh một ảnh đoán đúng của món bị đoán
nhầm sang (xanh).}{fig:img-pairs}

Năm cặp này chiếm 94 trên 825 lỗi. Mọi cặp đều là hai món cùng nhóm: bánh làm từ bột gạo hoặc món nước, đúng với những
nhóm món giống nhau đã thấy khi phân tích khám phá. Bánh canh và hủ tiếu bị nhầm theo cả hai chiều với tỉ lệ gần bằng
nhau.

\subsection{Mô hình sai ở đâu?}

\reportfig{fig12_error_analysis.png}{Tỉ lệ đoán sai trên test theo độ tin cậy của mô hình, theo việc ảnh có bản sao mang
nhãn khác, theo độ sáng và theo kích thước ảnh. Tỉ lệ sai chung là 16.4\%.}{fig:img-errors}

\begin{itemize}
  \item \textbf{Độ tin cậy} là dấu hiệu rõ nhất: ảnh có xác suất dưới 0.5 sai 61.5\% (576 ảnh), còn ảnh có xác suất trên
        0.95 chỉ sai 1.9\% (2 732 ảnh, hơn một nửa tập test).
  \item \textbf{Nhãn mâu thuẫn:} 37 ảnh test có bản sao mang nhãn khác bị đoán sai 51.4\%, so với 16.1\% ở các ảnh còn lại;
        với những ảnh này, nhãn trong dữ liệu có thể chính là chỗ sai.
  \item \textbf{Độ sáng và kích thước} gần như không ảnh hưởng: tỉ lệ sai từ 15.4\% đến 17.5\% giữa bốn nhóm độ sáng, và
        không có xu hướng rõ theo kích thước ảnh.
  \item \textbf{Số ảnh train:} hệ số tương quan Spearman giữa số ảnh train và F1 của từng món là -0.08 (p = 0.669). Món khó
        là món giống một món khác, không phải món ít ảnh.
\end{itemize}

\reportfig{fig13_confident_errors.jpg}{Mười sáu ảnh test bị đoán sai với độ tin cậy cao nhất (từ 0.99 trở lên); 3 ảnh
mang dấu \texttt{[trùng]}, tức có bản sao mang nhãn khác. Độ tin cậy trung vị là 0.54 với ảnh đoán sai và 0.98 với ảnh
đoán đúng.}{fig:img-confident}

\paragraph{Lỗi tự tin là lỗi của nhãn.} Khi mô hình rất chắc chắn mà vẫn sai, nhóm xem lại chính nhãn của ảnh. Trong 16
ảnh ở Hình~\ref{fig:img-confident}, ít nhất 9 ảnh khớp với món mô hình dự đoán hơn là nhãn gốc: hai ảnh gắn nhãn bánh
cuốn và nem chua là gỏi cuốn, ảnh gắn nhãn canh chua là cá kho tộ, ảnh gắn nhãn bún bò Huế là phở, ảnh gắn nhãn bún
riêu là bún bò Huế, ảnh gắn nhãn bánh đúc là bánh giò, ảnh gắn nhãn bánh bột lọc là bánh bèo, ảnh gắn nhãn bánh khọt là
bánh căn và ảnh gắn nhãn bánh tráng nướng là bánh cuốn. Chỉ 3 trong 16 ảnh có bản sao mang nhãn khác, nên phần lớn nhãn
sai này không tìm được bằng cách tìm ảnh trùng. Nhãn sai trong tập test làm điểm test thấp hơn chất lượng thật của mô
hình; nhóm không sửa nhãn để giữ nguyên tập test chính thức.

\section{Dự đoán ảnh mới}

Hàm \texttt{predict\_image()} chạy toàn bộ quy trình trên một file ảnh bất kỳ: đọc ảnh, đổi sang RGB, cùng phép biến đổi
kích thước và chuẩn hóa, rồi mô hình cuối cùng, và trả về 3 món có xác suất cao nhất. Khi xác suất của gợi ý đầu tiên
thấp (dưới 0.5), ứng dụng nên hiện cả 3 gợi ý thay vì một câu trả lời.

\reportfig{fig14_predictions.jpg}{Ba gợi ý của mô hình cho sáu ảnh test đọc từ file gốc, gồm bốn ảnh đoán đúng và hai ảnh
đoán sai (tên màu đỏ). Món đúng nằm trong ba gợi ý ở 5 trên 6 ảnh.}{fig:img-predict}

\section{Kết luận}

Mô hình cuối cùng, ResNet50 với khối cuối được fine-tune nhẹ, đạt accuracy 0.8363, macro-F1 0.8397 và top-3 0.9470
trên tập test, cao hơn bài báo gốc (accuracy 0.7754). Với cùng một kiểu đầu phân loại, đặc trưng học sẵn từ ImageNet là
bước tiến lớn nhất so với đặc trưng thủ công (macro-F1 từ 0.2400 lên 0.6945), và fine-tune nhẹ khối cuối cải thiện thêm
14.5 điểm. Lỗi tập trung ở các cặp món cùng nhóm và ở các ảnh bị gán nhãn sai trong dữ liệu.

\subsection*{Hạn chế}
\begin{itemize}
  \item Mỗi ảnh chỉ có một nhãn, trong khi một số ảnh có nhiều món; các biến thể vùng miền của một món được gộp chung.
  \item Nhãn sai vẫn còn trong val và test: 37 ảnh test có bản sao mang nhãn khác, và ít nhất 9 trong 16 lỗi tự tin
        nhất là ảnh bị gán nhãn sai. Nhóm không sửa nhãn, nên điểm test có thể thấp hơn chất lượng thật của mô hình.
  \item pHash chỉ tìm được bản sao của cùng một ảnh, không tìm được hai ảnh khác nhau chụp cùng một đĩa.
  \item Logistic Regression dùng tham số mặc định; fine-tune chỉ chạy 3 epoch trong khi validation vẫn đang tăng.
  \item Chỉ một mạng pretrained (ResNet50) được thử.
  \item Ảnh lấy từ Internet thường được chụp đẹp, khác ảnh chụp bằng điện thoại trong quán ăn.
\end{itemize}

\subsection*{Hướng mở rộng}
Fine-tune lâu hơn hoặc mở thêm \texttt{layer3}; thử các mạng khác như MobileNetV3 (nhẹ hơn, để chạy trên điện thoại) hoặc
ViT-B/16; kiểm tra lại nhãn của các ảnh mô hình đoán sai với độ tin cậy cao và các ảnh có bản sao mang nhãn khác, rồi
sửa hoặc bỏ chúng; tìm ảnh gần trùng bằng độ tương đồng
của đặc trưng CNN, ngoài pHash.

\begin{thebibliography}{9}
\raggedright
\bibitem{img-30vnfoods} T.-H. Do, D.-D.-A. Nguyen, H.-Q. Dang, H.-N. Nguyen, P.-P. Pham, D.-T. Nguyen.
  \textit{30VNFoods: A Dataset for Vietnamese Foods Recognition}. 2021 IEEE International Conference on Communication,
  Networks and Satellite (COMNETSAT), 2021. DOI: \url{10.1109/COMNETSAT53002.2021.9530774}.
\bibitem{img-resnet} K. He, X. Zhang, S. Ren, J. Sun. \textit{Deep Residual Learning for Image Recognition}. IEEE
  Conference on Computer Vision and Pattern Recognition (CVPR), 2016.
\end{thebibliography}
